\chapter{Introduction}
\label{ch:intro}
Radio astronomy studies the Universe through the radio waves emitted by
astronomical sources such as pulsars, galaxies and the neutral hydrogen gas
between them. By the time these signals reach the Earth they are extremely
weak, and radio telescopes are built to be very sensitive in order to detect
them. The same frequency bands, however, are also used by human technology:
mobile phone towers, satellites, aircraft and weather radar, television and
radio broadcasts, and electrical equipment such as power lines. Signals from
these sources leak into the telescope and are usually many orders of
magnitude stronger than the astronomical signal. This unwanted signal is known
as radio frequency interference (RFI).

RFI is different from the ordinary noise that is present in every radio
observation. The thermal noise of the receiver and the sky is random and close
to Gaussian, and its statistics are well understood. RFI on the other hand is
strong, has structure in both time and frequency, and changes during the
observation. Depending on its source it can appear as a narrow line at a single
frequency that lasts for the whole observation, a short burst that covers
almost all frequencies at once, or a broad patch that appears and fades within
a few minutes. Before any scientific analysis, the parts of the data affected
by RFI have to be found and removed, a step called RFI flagging. Modern
telescopes produce far too much data for this to be done by hand, so flagging
has to be automatic.

Radio data can be arranged as a two-dimensional image with frequency along one
axis and time along the other, called a spectrogram. In this picture RFI shows
up as lines, bands and patches, and flagging means deciding for every pixel
whether it is affected by RFI or not. This is the same problem that computer
vision calls semantic segmentation, and it can therefore be approached with the
convolutional neural networks used for image segmentation. The
U-Net\footcite{ronneberger2015} is the standard network for this task, and
Akeret et al.\footcite{akeret2017} were the first to apply it to RFI, in a code
called \tfunet{}. A different approach was proposed by Chen et
al.\footcite{chen2023}, who first label pixels with a statistical threshold and
then refine these labels with a conditional random field (CRF) that takes the
neighbourhood of each pixel into account.

In this project we treat RFI flagging as an image segmentation problem. We
reproduce the U-Net method of Akeret et al., design an improved network called
\hybrid{}, and test both on synthetic data and on real observations from the
Low-Frequency Array (LOFAR) that were labelled by a human expert.

\section{Motivation and Objectives of this Work}
Pulsars are a good example of why RFI flagging matters. Pulsars are rotating
neutron stars that emit very regular pulses, and millisecond pulsars are
regular enough to be used as precise clocks. Pulsar timing arrays time such
pulsars to search for low-frequency gravitational waves, which needs the
arrival times of the pulses to be measured to better than a microsecond. RFI
that is not removed distorts the measured pulse profiles and therefore the
arrival times. The tools commonly used to clean such data, for example
\texttt{paz} from PSRCHIVE\footcite{vanstraten2012psrchive} and
\texttt{clfd}\footcite{clfd}, rely on simple statistics such as median
filters. Since RFI is not Gaussian, these tools can miss weak RFI or, when made
more aggressive, remove parts of the real signal.

Neural networks can learn what RFI looks like directly from examples instead
of relying on a fixed statistical rule. Most of these networks, however, are
developed and tested on synthetic data, where the exact position of every RFI
pixel is known. Real observations are much harder: RFI is rarer, a large part
of it is faint, and the labels available for training are produced by other
automatic flaggers and are themselves often wrong. It is therefore not clear
how much of the reported performance carries over to real data, or which parts
of a network are actually responsible for it.

This project approaches RFI detection as a pixel-wise segmentation problem on
spectrograms. We use a synthetic dataset that we generated ourselves, a
simulated dataset from the Hydrogen Epoch of Reionization Array (HERA), and
real LOFAR data whose test images were labelled by a human expert. Our specific
objectives are:
\begin{itemize}
  \item To reproduce the U-Net method of Akeret et al.\ (\tfunet{}) on our own
    data, understand why it initially performed poorly, and correct it.
  \item To design an improved network for RFI detection (\hybrid{}) and compare
    it with \tfunet{} and with published methods.
  \item To test both networks on real telescope data labelled by a human
    expert, and not only on synthetic data.
  \item To find out which parts of the improved network actually make it
    better, including whether a CRF refinement step helps.
  \item To study whether a small number of expert labels can improve the
    network further, and how these labels should be chosen.
\end{itemize}

The poor initial performance of \tfunet{} was traced back to its training
settings, and the improved network was compared with published methods on the
109 expert-labelled LOFAR images. We also removed the components of the network
one at a time to measure what each contributes, and fine-tuned it using expert
labels. Through this analysis, we aim to understand not only how well a neural
network can flag RFI in real observations, but also why it works.
